\documentclass[10pt,english]{article}

\usepackage{amsfonts,url,tikz,amsmath}

% Paper setup
\evensidemargin=0in
\oddsidemargin=0in
\textwidth=6.25in
\topmargin=-0.5in
\headheight=0.0in
\headsep=0.5in
\textheight=9.0in
\footskip=0.5in

\input{institution}

\newcommand{\Interior}[1]{\ensuremath{{#1}^{\circ}}}
\newcommand{\Closure}[1]{\ensuremath{\overline{#1}}}
\newcommand{\Complement}[1]{\ensuremath{{#1}^{c}}}

\newcommand{\Expect}{\ensuremath{\mathrm{E}}}
\newcommand{\vecnot}{\underline}
\newcommand{\RealNumbers}{\ensuremath{\mathbb{R}}}
\newcommand{\RationalNumbers}{\mathbb{Q}}
\newcommand{\ComplexNumbers}{\mathbb{C}}
\newcommand{\Real}{\mathrm{Re}}
\newcommand{\Span}{\mathrm{span}}
\newcommand{\Rank}{\mathrm{rank}}
\newcommand{\Nullity}{\mathrm{nullity}}
\newcommand{\Trace}{\mathrm{tr}}
\newcommand{\Diag}{\mathrm{diag}}
\newcommand{\dd}{\mathrm{d}}
\DeclareMathOperator*{\esssup}{ess\,sup}

% Use < , > inner product
\newcommand{\inner}[2]{{\left\langle #1 \mskip2mu , #2 \right\rangle}}
\newcommand{\tinner}[2]{{\langle #1 \mskip1mu , #2 \rangle}}

% Use < | > inner product
%\newcommand{\inner}[2]{{\left\langle #1 \mskip2mu \middle| \mskip2mu #2 \right\rangle}}
%\newcommand{\tinner}[2]{{\langle #1 \mskip1mu | \mskip1mu  #2 \rangle}}



\begin{document}

\title{ECE 586: Vector Space Methods \\ Lecture 21: Frequently Asked Questions}
\author{Henry D. Pfister \\ \univ}
\date{October 18, 2026}

\maketitle

\section{Convex Sets and Convex Functions}

\paragraph{How is a convex set different from a subspace?}
A convex set contains every line segment joining two of its points.
A subspace is closed under all linear combinations, so it is convex and contains zero.
A ball or a shifted half space can be convex without being a subspace (slide 1, ``Convexity'').

\paragraph{Can a convex function have several minimizers?}
Yes. A constant function has every feasible point as a minimizer.
Convexity guarantees that every local minimum on a convex feasible set is global, not that its location is unique.
Strict convexity gives at most one minimizer, if a minimizer exists (slide 2, ``Convex Optimization'').

\paragraph{Does strict convexity guarantee that a minimum is achieved?}
No. For example, $f(x)=e^x$ is strictly convex on $\mathbb{R}$ but has infimum zero and no minimizer.
Existence requires an additional argument, such as continuity on a nonempty compact feasible set.
For the squared distance to a nonempty closed convex set in a Hilbert space, the projection theorem supplies existence without requiring compactness.

\paragraph{Why does a tangent plane give a global lower bound for a convex function?}
For a differentiable convex function, slide 3, ``Convex Optimization and Derivatives,'' gives
\[
f(x)\geq f(x_0)+\inner{x-x_0}{\nabla f(x_0)}.
\]
Convexity controls the function along the segment from $x_0$ to $x$; taking the initial slope gives this bound.
A zero gradient therefore proves global optimality for an unconstrained problem.

\paragraph{At a constrained minimum, must the gradient be zero?}
No. At a boundary point, the negative gradient may point outside the feasible set.
For example, minimizing $f(x)=x$ over $[0,\infty)$ gives $x^*=0$ even though $f'(0)=1$.
For a differentiable convex function on a convex set $A$, the appropriate condition is $\inner{x-x^*}{\nabla f(x^*)}\geq0$ for every $x\in A$.

\section{Projection onto Convex Sets}

\paragraph{Why do we require the set to be nonempty, closed, and convex?}
Nonemptiness provides feasible candidates, closedness prevents the nearest point from being excluded, and convexity gives uniqueness in a Hilbert space.
For example, projecting $0$ onto $(0,1)$ has no solution, while projecting $0$ onto $\{-1,1\}$ has two nearest points.
Slide 4, ``Projection onto Convex Sets,'' should be read with the nonempty-set assumption.

\paragraph{What does the projection inequality mean geometrically?}
For a real Hilbert space, $p=P_A(v)$ is characterized by
\[
\inner{v-p}{u-p}\leq0\quad\text{for every }u\in A.
\]
Thus the residual from $p$ toward $v$ makes a nonacute angle with every feasible direction from $p$ into the set.
In a complex Hilbert space, use the real part of the inner product in this inequality.

\paragraph{How does projection onto a subspace fit this inequality?}
For a subspace, both $p+tw$ and $p-tw$ are feasible whenever $w$ is in the subspace.
Applying the inequality with both signs forces the residual to be orthogonal to every $w$.
This recovers the orthogonality condition used in Videos 17 and 18.

\paragraph{Is projection onto a convex set a linear map?}
Generally no. Projection onto $[0,\infty)$ in $\mathbb{R}$ is $P_A(x)=\max\{x,0\}$, which is nonlinear.
It still fixes points already in the set and is nonexpansive: $\|P_A(v)-P_A(w)\|\leq\|v-w\|$.
For a closed subspace, the orthogonal projection is linear.

\end{document}
